\section{引言：课程概览与学习目标}

CS224N（Natural Language Processing with Deep Learning）是 Stanford 大学最受欢迎的 NLP 课程之一，由 Christopher Manning 教授主讲。本节课作为第一讲，涵盖三大主题：（1）课程介绍与学习目标；（2）人类语言与词义表示；（3）Word2Vec 算法详解。

\begin{figure}[H]
\centering
\includegraphics[width=0.95\textwidth]{slides-images/slide-000.jpg}
\caption{CS224N/Ling284 课程封面：Natural Language Processing with Deep Learning\protect\footnotemark}
\end{figure}
\footnotetext{Slides 第1页。}

\subsection{课程学习目标}

Manning 教授明确列出了三个核心学习目标：

\begin{importantbox}{CS224N 三大学习目标}
\begin{enumerate}
    \item \textbf{掌握深度学习在 NLP 中的基础与前沿方法}：从词向量、前馈神经网络、循环网络、注意力机制，到 Transformer、编码器-解码器模型、大语言模型的预训练与后训练、模型适配、可解释性、智能体等
    \item \textbf{理解人类语言的特性与计算挑战}：了解语言结构的复杂性，以及为什么让计算机理解和生成人类语言如此困难
    \item \textbf{具备构建实际 NLP 系统的能力}：不仅学习理论，还能在实际工作中独立搭建文本分类、信息抽取、机器翻译等系统
\end{enumerate}
\end{importantbox}

\begin{figure}[H]
\centering
\includegraphics[width=0.95\textwidth]{slides-images/slide-003.jpg}
\caption{课程学习目标：基础与前沿方法、语言理解、系统构建三个维度\protect\footnotemark}
\end{figure}
\footnotetext{Slides 第4页。}

\subsection{课程结构与作业安排}

课程包含四个编程作业和一个期末项目：

\begin{itemize}
    \item \textbf{Assignment 1}：入门作业（Jupyter Notebook），探索词向量
    \item \textbf{Assignment 2}：数学推导 + PyTorch 入门，构建依存句法分析器
    \item \textbf{Assignment 3--4}：使用 PyTorch + GPU 完成机器翻译和 Transformer 相关项目
    \item \textbf{期末项目}：默认项目（有脚手架）或自选项目（1--3人组队）
\end{itemize}

\begin{figure}[H]
\centering
\includegraphics[width=0.95\textwidth]{slides-images/slide-005.jpg}
\caption{评分政策：作业占 48\%，期末项目占 49\%，课堂参与占 3\%\protect\footnotemark}
\end{figure}
\footnotetext{Slides 第6页。}

\subsection{本章小结}

CS224N 是一门从底层到前沿的 NLP 课程，强调理论与实践并重。课程从词向量出发，逐步深入到 Transformer、大语言模型等前沿技术，最终要求学生具备独立构建 NLP 系统的能力。

%% ========== Part 2: 人类语言与词义 ==========

